\begin{examples}[7.3.8]
\label{IV.2.s7.3.8-fr}
Las propiedades $\boldsymbol P(Z,k)$ siguientes
satisfacen las condiciones $(\mathrm P_{\mathrm I})$,
$(\mathrm P_{\mathrm{II}})$ y $(\mathrm P_{\mathrm{III}})$:
\begin{enumerate}
\item[(i)] (también denotada $(\mathrm i_n)$) $Z$ tiene coprofundidad
$\leqslant n$.
\item[(ii)] $Z$ es un preesquema de Cohen-Macaulay.
\item[(iii)] (también denotada $(\mathrm{iii}_n)$) $Z$ satisface
$(\mathrm S_n)$.
\item[(iv)] $Z$ es regular.
\item[(v)] (también denotada $(\mathrm v_n)$) $Z$ satisface
$(\mathrm R_n)$.
\item[(vi)] $Z$ es reducido.
\item[(vii)] $Z$ es normal.
\end{enumerate}
Para las propiedades (ii) a (vii), esto se deduce de (6.6.1), que de
hecho demuestra la condición más fuerte $(\mathrm P'_{\mathrm I})$.
Para (i), la propiedad $(\mathrm P_{\mathrm{II}})$ se deduce de
(6.6.2); la propiedad $(\mathrm P_{\mathrm I})$ se deduce, por el
mismo razonamiento que en (6.6.1, (i)), de (6.3.2) y del hecho de que
un preesquema regular tiene coprofundidad 0.

Además, de (6.7.1) se deduce que las propiedades (i), (ii) y (iii) son
\emph{geométricas}; en virtud de (6.7.8), lo mismo ocurre con las siguientes:
\begin{enumerate}
\item[(iv$'$)] $Z$ es geométricamente regular.
\item[(v$'$)] (también denotada $(\mathrm v'_n)$) $Z$ posee la
propiedad $(\mathrm R_n)$ geométrica.
\item[(vi$'$)] $Z$ es geométricamente reducido.
\item[(vii$'$)] $Z$ es geométricamente normal.
\end{enumerate}
\end{examples}
